% !TeX root = ../../Book4.tex
% 纲要标题：### 17.7*　有理 Hodge 类与线丛
% 纲要位置：几何学计划纲要.md，第 4179 行。
% 正文状态：未撰写；本文件目前是框架。
% 写作范围、先修与边界：先修为 IV.16 的全纯线丛和层上同调、IV.17.4 的 Hodge 分解以及第二卷的整上同调与系数变换；射影代数圈的解释另用 IV.20 的比较，不要求先完成 V.27。先在紧 Kähler 流形上由指数正合列证明 Lefschetz (1,1) 定理，核对整系数、挠元和第一 Chern 类的像；有理 Hodge 猜想只在光滑复射影簇上准确陈述，代数圈类的更完整形式指向 V.27。明确有理与整数版本、除子与高余维的区别，不承担 Hodge 猜想的一般证明。

% 纲要内容开始
% 
% > 写作范围与外部输入：先修为 IV.16 的全纯线丛和层上同调、IV.17.4 的 Hodge 分解以及第二卷的整上同调与系数变换；射影代数圈的解释另用 IV.20 的比较，不要求先完成 V.27。先在紧 Kähler 流形上由指数正合列证明 Lefschetz (1,1) 定理，核对整系数、挠元和第一 Chern 类的像；有理 Hodge 猜想只在光滑复射影簇上准确陈述，代数圈类的更完整形式指向 V.27。明确有理与整数版本、除子与高余维的区别，不承担 Hodge 猜想的一般证明。
% 
% * 有理 Hodge 类与整上同调的区别
% * 指数正合列与线丛的第一 Chern 类
% * 线丛的整 Hodge 类与 Lefschetz 定理
% * 代数圈的上同调类与有理 Hodge 猜想
% * 除子与高余维问题的区别
% 
% ---
% 
% 纲要内容结束

\OptionalSection{有理 Hodge 类与线丛}{g4:sec:1707}

待撰写

\subsection{有理 Hodge 类与整上同调的区别}
\label{g4:subsec:1707:01}
待撰写

\subsection{指数正合列与线丛的第一 Chern 类}
\label{g4:subsec:1707:02}
待撰写

\subsection{线丛的整 Hodge 类与 Lefschetz 定理}
\label{g4:subsec:1707:03}
待撰写

\subsection{代数圈的上同调类与有理 Hodge 猜想}
\label{g4:subsec:1707:04}
待撰写

\subsection{除子与高余维问题的区别}
\label{g4:subsec:1707:05}
待撰写

% 纲要条目排为编号小节标题，并进入目录与书签；说明段仍保留于纲要注释中。
